
%\needspace{4\baselineskip}
\color{uniblue}{\subsection{Восстановление нормального режима (ВНР)}}
\color{black}

\begin{enumerate}[label=\arabic{section}.\arabic{subsection}.\arabic*, labelsep=4pt, leftmargin=0pt, itemindent=57pt]

\item
Функция ВНР предназначена для автоматического восстановления схемы доаварийного режима после действия АВР. Устройство может выполнять обратное переключение, что  подразумевает отключение секционного выключателя (выключателя резервного ввода) и включение основного рабочего ввода.

\item
В алгоритме ВНР предусмотрено два режима работы: без перерыва питания потребителей (ВВ-СВ) и с перерывом питания (СВ-ВВ). В первом случае работа осуществляется без перерыва питания, включается 
основной рабочий ввод и потом отключается секционный выключатель (выключатель резервного ввода). Во  втором  случае  сначала  отключается  секционный  выключатель  затем  проверяется  его  включенное состояние и после этого включается выключатель основного рабочего ввода. Данные режимы выбираются уставкой SGF2 <<Очередность\_ВНР>>. 

\item
Критерием пуска алгоритма ВНР служит появление рабочего напряжения на ТН до ввода. Алгоритм ВНР может работать только после срабатывания АВР. При оперативном управлении выключателем алгоритм ВНР не запустится. Действие функции ВНР блокируется, если выведена (заблокирована) функция АВР. 

\item
Готовность  схемы  ВНР  и  активация  соответствующего  сигнала <<ВНР:  Ожидание  ВНР>> осуществляется при аварийном отключении основного рабочего ввода (активный сигнал <<В: Отключен>>). 
Обязательным  условием  является  отсутствие  напряжения  до  вводного  выключателя,  а  также  наличие разрешения от АВР <<АВР: Вкл. рез.ВВ или СВ>>. Временным таймером Т5 продлевается сигнал разрешения АВР по спаду сигнала на 10 мс. При выполнении всех перечисленных условий формируется выходной сигнал  ожидания  ВНР,  который  сохраняется  в  энергонезависимой  памяти  до  появления  сигналов  для сброса триггера.

\item
Описание алгоритмов работы ВНР:
\begin{enumerate}[label=\arabic*), noitemsep]
    \item Алгоритм работы схемы ВВ-СВ:
    \begin{itemize}
        \item При наличии условий для формирования выходного сигнала <<ВНР: Ожидание ВНР>> и сохранении напряжения на ТН до ввода, собираются условия для включения рабочего ВВ. Включение выключателя  может  выполняться  как  с  контролем  синхронизма,  так  и  без  контроля.  Контроль  синхронизма  вводится уставкой  SGF4.  Без  контроля  синхронизма  команда  на  включение  основного  рабочего  вводного выключателя  формируется  по  истечению  регулируемой  выдержки  времени  таймера  Т2.  Обязательное условие для включения основного ВВ - отсутствие сигнала блокировки включения из логики КСВ.
        \item При сохранении условий на включение основного рабочего вводного выключателя, истечению выдержки времени таймера Т2 и  включенного положения рабочего ВВ (активный сигнал <<В: Включен>>) формируются условия для отключения СВ (выключателя резервного ввода). Таймер Т6 кратковременно пропускает  сигнал  на  время  5  мс  в сторону триггера. С выхода триггера  активируется таймер Т3, по  истечении выдержки времени которого формируется  выходная  команда  в  терминал устройства  на  отключение  СВ  (выключатель  резервного  ввода). Совместно с командой отключения выключателя выполняется сброс триггеров.         
    \end{itemize}
    \item Алгоритм работы схемы СВ-ВВ:
    \begin{itemize}
        \item При наличии условий ожидания ВНР и напряжения на ТН до вводного выключателя, по истечении выдержки времени таймера Т2 формируется команда в терминал СВ (выключателя резервного ввода) для отключения  выключателя.  После  отключения  СВ  (выключателя  резервного  ввода)  на  дискретный  вход терминала  основного  рабочего  ввода  поступает  дискретный  сигнал  отключенного  положения  СВ (выключателя резервного ввода) <<Рез.ВВ/СВ отключен>>. С выхода триггера  активируется  таймер  Т3,  по  истечении выдержки времени которого формируется  выходная  команда в привод на включение рабочего вводного выключателя с контролем или без контроля отсутствия напряжения на шинах (уставка  SGF3 <<Контроль\_U>>).  Обязательное условие для  включения  основного  ВВ  --  отсутствие  сигнала  блокировки включения из логики КСВ. Совместно с командой включения выключателя выполняется сброс триггеров, возвращая схему в исходное состояние.
    \end{itemize}
\end{enumerate}

\item
Сброс триггеров осуществляется при выполнении следующих условий:
\begin{itemize}
    \item функция ВНР выведена из работы; 
    \item задержка включения выключателя с КС <<ВНР: Задержка вкл. с КС>>; 
    \item активный сигнал из логики АВР <<АВР: Вкл. рез.ВВ или СВ>>; 
    \item активный сигнал из логики запрета АВР <<ЗАВР: Запрет АВР>>; 
    \item блокировка включения выключателя из логики КСВ; 
    \item активный  внешний  дискретный  сигнал  блокировки  ВНР  <<Блокировка  ВНР>>  (на  усмотрение Заказчика); 
    \item отсутствие сигнала РФК из логики КСВ. 
\end{itemize}

\item
Ввод ВНР в работу на этапе параметрирования устройства осуществляется программным переключателем SGF1 <<Ввод\_функции>>, при этом необходимо введенное состояние функции АВР (активный сигнал <<АВР: Ввод>>). Информация о введенной функции отображается выходным сигналом <<ВНР: Ввод>>.

\item
Функция <<ВНР>> может быть выведена из работы оперативно путем активации сигнала <<ОВ ВНР>>. Выведенное состояние функции характеризуется активным состоянием сигнала <<ВНР: Оперативный вывод>>.

\item
Алгоритм работы функции <<ВНР>> выполнен в соответствии с рисунком \ref{vnr:img1}. В таблице \ref{vnr:tbl1} приведены параметры, необходимые для настройки функции <<ВНР>>.

\vspace{3mm}
\begin{figure}[H]
\centering
\includegraphics[width=1\textwidth,height=1\textheight,keepaspectratio]{\fbpath/310. ВНР/_latex/img/vnr.pdf}
\captionof{figure}{Функциональная схема <<ВНР>>}\label{vnr:img1}
\end{figure}

\small
\begin{longtable}{|>{\centering\arraybackslash}m{5.3cm}|>{\centering\arraybackslash}m{3.3cm}|>{\centering\arraybackslash}m{4.2cm}|>{\centering\arraybackslash}m{1.8cm}|>{\centering\arraybackslash}m{1cm}|}
\caption{Параметры для настройки функции <<ВНР>>\hfill\vspace{-0.5\baselineskip}}\label{vnr:tbl1}\\ 
\hline
\rowcolor{gray!30}
Параметр (Параметр на ИЧМ) & Условное обозначение на схеме & Значение/ Диапазон & Единица измерения & Шаг \\ 
\hline
\endfirsthead
\caption*{\hspace{3pt}\emph{Продолжение таблицы \ref{vnr:tbl1}\hfill\vspace{-0.5\baselineskip}}} \\ % сделано по ГОСТ 2.105 п.6.8.7
\hline
\rowcolor{gray!30}
Параметр (Параметр на ИЧМ) & Условное обозначение на схеме & Значение/ Диапазон & Единица измерения & Шаг \\ 
%\hline
\endhead
%\hline
\endfoot
%\hline
\endlastfoot
%==+t1*ANSR1|VV_LVNSR
\centering Ввод функции в работу (Ввод\_функции) & \centering SGF1 & \centering 0 = Не предусмотрено\\1 = Предусмотрено & \centering -- & \centering \arraybackslash -- \\
\hline
\centering Порядок действия (Очередность\_ВНР) & \centering SGF2 & \centering 1 = ВВ - СВ\\2 = СВ - ВВ & \centering -- & \centering \arraybackslash -- \\
\hline
\centering Контроль отсутствия напряжения (Контроль\_U) & \centering SGF3 & \centering 0 = Не предусмотрено\\1 = Предусмотрено & \centering -- & \centering \arraybackslash -- \\
\hline
\centering Контроль синхронизма (Контр\_синхр) & \centering SGF4 & \centering 0 = Не предусмотрено\\1 = Предусмотрено & \centering -- & \centering \arraybackslash -- \\
\hline
\centering Время ввода ВНР (Tввода) & \centering T1 & \centering 0,00 ... 300,00 & \centering с & \centering \arraybackslash 0,01 \\
\hline
\centering Выдержка времени 1 операции (Tопер1) & \centering T2 & \centering 0,00 ... 300,00 & \centering с & \centering \arraybackslash 0,01 \\
\hline
\centering Выдержка времени 2 операции (Tопер2) & \centering T3 & \centering 0,00 ... 50,00 & \centering с & \centering \arraybackslash 0,01 \\
\hline
\centering Длительность задержки включения с контролем синхронизма (Тзад\_вкл) & \centering T4 & \centering 0,00 ... 300,00 & \centering с & \centering \arraybackslash 0,01 \\
\hline
%===t1
\end{longtable}
\normalsize

\end{enumerate}
